%! TeX program = lualatex
\documentclass[../main.tex]{subfiles}
\begin{document} \section{Problem-solving practice}

\begin{exercise}
  Evaluate \(\int 4u^{3} (1 + \sqrt{u}) \;du\).

  \begin{enumerate}[label=\Alph*.]
    \item \( u^{4} \left(u + \frac{2}{3} u^{3/2}\right) + C\)
    \item \( u^{4} \left(u + u^{3/2}\right) + C\)
    \item \(u^{4} + \frac{2}{9}u^{9/2} + C\)
    \item \(u^{4} + \frac{8}{9}u^{9/2} + C\)
    \item \(4u^{4} + u^{9/2} + C\)
  \end{enumerate}
\end{exercise}

\begin{exercise}
  Which expression defines \(\int_{0}^{3} x^{3} - 2x \;dx\)?

  \begin{enumerate}[label=\Alph*.]
    \item \(\lim_{n \to \infty} \sum_{i=1}^{n} \left( \frac{3i^{3}}{n} - \frac{6i}{n} \right)\).
    \item \(\lim_{n \to \infty} \sum_{i=1}^{n} \left( \frac{9i}{n^{2}} - \frac{6i}{n^{3}} \right)\).
    \item \(\lim_{n \to \infty} \sum_{i=1}^{n} \left( \frac{9i}{n^{2}} - \frac{27i}{n^{3}} \right)\).
    \item \(\lim_{n \to \infty} \sum_{i=1}^{n} \frac{3}{n}(i^{3} - 2i)\).
    \item \(\lim_{n \to \infty} \sum_{i=1}^{n} \left( \frac{i}{n^{2}} - \frac{2i}{n} \right)\).
    \item None of the above.
  \end{enumerate}
\end{exercise}

\begin{exercise}
  Suppose \(F'(x) = \frac{x}{\sqrt{1+x^{2}}}\) and \(F(1) = 0\). Find \(F(x)\).

  \begin{enumerate}[label=\Alph*.]
    \item \(F(x) = \sqrt{\arctan(x)}\).
    \item \(F(x) = \sqrt{\arctan(x)} - \sqrt{\pi/4}\).
    \item \(F(x) = \sqrt{x^{2} + 1} - \sqrt{2}\).
    \item \(F(x) = \sqrt{x^{2} + 1} + \sqrt{2}\).
    \item \(F(x) = \sqrt{x^{2} + 1}\).
  \end{enumerate}
\end{exercise}

\begin{exercise}
  Evaluate \(\lim_{x \to \infty} \frac{e^{x} + \ln(x)}{5^{x}}\).

  \begin{enumerate}[label=\Alph*.]
    \item \(y = x - 3\)
  \end{enumerate}
\end{exercise}

\begin{exercise}
  Find the area between the curves \(y = \sqrt{x}\) and \(y = x^{5}\).

  \begin{enumerate}[label=\Alph*.]
    \item 
  \end{enumerate}
\end{exercise}

\begin{exercise}
  Suppose \(f'(x) = \int_{0}^{x} \sin(t^{2}) \;dt\).  Which of the following statements is true?

  \begin{enumerate}[label=\Alph*.]
    \item \(f(x)\) is decreasing over the interval \([-\sqrt{\pi/2}, 0]\).
    \item \(f(x)\) is not defined for \(x < 0\).
    \item \(f(x)\) is concave down over the interval \([0, \sqrt{\pi/2}]\).
    \item \(f(x)\) is concave up over the interval \([0, \sqrt{\pi/2}]\).
    \item \(f(x)\) has an inflection point when \(x = \sqrt{\pi/2}\).
  \end{enumerate}
\end{exercise}

\begin{exercise}
  Evaluate \(\lim_{x \to \infty} \frac{\ln(2x + 1)}{3\ln(x)-1}\).

  \begin{enumerate}[label=\Alph*.]
    \item \(0\).
    \item \(1/3\).
    \item \(2/3\).
    \item \(1\).
    \item \(\infty\).
  \end{enumerate}
\end{exercise}

\begin{exercise}
  Evaluate \(\int_{2}^{1} \sec^{2}\left(\frac{1}{\sqrt[3]{x}}\right) \;dx\).

  \begin{enumerate}[label=\Alph*.]
    \item \(0\)
    \item \(-\tan(\ln(2))\)
    \item The integral cannot be evaluated.
  \end{enumerate}
\end{exercise}

\begin{exercise}
  Find the left-endpoint approximation of \(\int_{2}^{6} \ln(x) \;dx\) using \(4\) rectangles of equal width.

  \begin{enumerate}[label=\Alph*.]
    \item \(\ln(2)\ln(3)\ln(4)\ln(5)\)
    \item \(\ln(3)\ln(4)\ln(5)\ln(6)\)
    \item \(\ln(720)\)
    \item \(\ln(120)\)
    \item \(\ln(360)\)
  \end{enumerate}
\end{exercise}

\begin{exercise}
  Find absolute extrema of \(h(x) = 2x^{3} + 3x^{2} - 36x - 2\) over the interval \([-3,2]\).
\end{exercise}

\begin{exercise}
  A window has the shape of a rectangle topped by a semicircle. Its perimetre is \(100\) metres.  Maximize the area of the window.
\end{exercise}
\end{document}
